\chapter*{Acknowledgements}
\addcontentsline{toc}{chapter}{Acknowledgements}

\noindent
\textbf{English version}

\noindent
After two theses, I have learned the hard way that writing exhaustive acknowledgements is like searching for a proof: it exists somewhere, but nobody promises it can be found efficiently. In these pages I have defined, assumed, stated lemmas and, probably more than once, contradicted myself. Yet this is the only section in which I can state something without a single quantifier: I did not get here alone.

To my advisors, Nicola and Massimo. Thank you for your patience, your rigor, and the rare ability to say \curlyquotes{this step is wrong} or even \curlyquotes{I can't make sense of any of this}. You taught me that a good question is worth more than a quick answer, and that a good answer is worth checking twice. You are not only mentors but also a source of aspiration. Whatever is correct in these pages I owe largely to you. Whatever is not is, as tradition demands, left as an exercise for the reader.

To the thesis reviewers and the graduation committee. Thank you for giving your time to a text whose central message is that finding short proofs is hard, and for then asking you to read several that are extremely long. I hope the effort feels worthwhile, and I promise that the next proof will be shorter, although I cannot prove it without breaking some assumption.

To my friends and department colleagues, Sara, Alessio, Armando, Ivan and Ionut, for the whiteboards full of lemmas that did not survive, for the coffee that did, and for the conversations in which \curlyquotes{it is obviously true} turned into \curlyquotes{it is obviously false} within ten minutes. A thesis is a solitary object, but the years of study behind it were not.

To the people who, unfortunately, have to deal with me outside academia. Thank you for never asking what a proof system is and, if only for that, for preserving my sanity. For the evenings and the discussions in which the only open problem was where to go out. And yes, I can finally answer the question \curlyquotes{so, when do you graduate?} with a date. Thank you to Sara, Gaia and Clara who, unfortunately for them, have been with me forever, and to Lorenzo, Micol, Giulia and Valentina for putting up with me every day, provided that I actually get in touch, of course.

To my family. Thank you for your support, for your sacrifices, and for believing in me long before there was anything to believe in. You gave me the freedom to choose and the certainty of always having a place to come back to. I hope this milestone, like the future ones, feels like yours too, because it is: every line of this thesis has a piece of your work behind it, and no result of mine will ever be as important as this one.

To Camilla. Our relationship developed just like the search for a result: the outcome was a foregone conclusion, but it was necessary to detour and explore other paths, some of them difficult, in order to obtain it. You listened to explanations you were under no obligation to follow, celebrated results you could not decipher, and consoled failures you could not verify. I always tell you this and I claim it here: you are the only person I consider above myself. You are out of this world; indeed, I will not deny that the green dress you wore the first time we met suggested you were an alien. Luckily, my obsession is the study of complex, intractable and incomprehensible things.

To myself, for not abandoning the project during the weeks when the main result seemed impossible to reach, and for finally understanding that a deadline is a theorem you only believe once you have proved it.

Finally, to everyone I forgot to name: the omission is a limit of my memory, but not of my gratitude. To ensure correctness I had to sacrifice completeness, but I cannot even prove my own correctness. Damned Gödel and his two stupid theorems.

\noindent
\textit{Simone Bianco, Rome, October 2026}

\quad

\noindent
\textbf{Versione italiana}

\noindent
Dopo ben due tesi ho imparato, a mie spese, che scrivere dei ringraziamenti esaustivi è come cercare una dimostrazione: esiste da qualche parte, ma nessuno promette che si possa trovare in modo efficiente. In queste pagine ho definito, assunto, enunciato lemmi e, probabilmente più di una volta, contraddetto me stesso. Tuttavia, questa è l'unica sezione in cui posso affermare qualcosa senza un solo quantificatore: non sono arrivato fin qui da solo.

Ai miei relatori, Nicola e Massimo. Grazie per la pazienza, il rigore e la rara capacità di dire \curlyquotes{questo passaggio è sbagliato} o persino \curlyquotes{non si capisce niente}. Mi avete insegnato che una buona domanda vale più di una risposta rapida, e che una buona risposta vale la pena di essere controllata due volte. Non siete solo mentori ma anche fonte di aspirazione. Tutto ciò che è corretto in queste pagine lo devo in larga parte a voi. Tutto ciò che non lo è, come vuole la tradizione, è lasciato per esercizio al lettore.

Ai revisori della tesi e alla commissione di laurea. Grazie per aver dedicato il vostro tempo a un testo il cui messaggio centrale è che trovare dimostrazioni brevi è difficile, e per avervi poi chiesto di leggerne svariate che sono lunghissime. Spero che lo sforzo vi sembri valso la pena e prometto che la prossima dimostrazione sarà più breve, anche se non sono in grado di dimostrarlo a meno di rompere qualche assunzione.

Ai miei amici e colleghi di dipartimento, Sara, Alessio, Armando, Ivan e Ionut, per le lavagne piene di lemmi che non sono sopravvissuti, per il caffè che invece ce l'ha fatta e per le conversazioni in cui la frase \curlyquotes{è ovviamente vero} diventava \curlyquotes{è ovviamente falso} nel giro di dieci minuti. Una tesi è un oggetto solitario, ma gli di studio anni che ci sono dietro non lo sono stati.

Alle persone che, purtroppo per loro, hanno a che fare con me quotidianamente. Grazie per non aver mai chiesto cosa sia un sistema di dimostrazione e che, anche solo per questo, hanno preservato la mia sanità mentale. Per le serate e le discussioni in cui l'unico problema aperto era dove uscire. E sì, finalmente posso rispondere alla domanda \curlyquotes{allora, quando ti laurei?} con una data. Grazie a Sara, Gaia e Clara che mi accompagnano da sempre e grazie a Lorenzo, Micol, Giulia e Valentina per sopportarmi ogni giorno, a patto che io abbia la decenza di farmi vivo ovviamente.

Alla mia famiglia. Grazie per il sostegno, per i sacrifici e per aver creduto in me molto prima che ci fosse qualcosa in cui credere. Mi avete lasciato la libertà di scegliere e la certezza di avere sempre un posto in cui tornare. Spero che questo traguardo, come quelli futuri, sia percepito da voi anche come vostro, perché lo è: ogni riga di questa tesi ha dietro un pezzo del vostro lavoro e nessun risultato mio sarà mai importante quanto questo.

A Camilla. Il nostro rapporto si è sviluppato proprio come la ricerca di un risultato: l'esito sembrava vero fin da subito, ma è stato necessario sviare ed esplorare altre strade, alcune difficili, per poterlo dimostrare. Hai ascoltato spiegazioni che non eri tenuta a seguire, festeggiato risultati che non potevi decifrare e consolato fallimenti che non potevi verificare. Te lo dico sempre e qui lo rivendico: sei l'unica persona che ritengo al di sopra di me. Sei al di fuori di questo mondo, difatti non nego che il vestito verde da te indossato al nostro primo incontro suggerisse il tuo essere aliena. Fortuna vuole che la mia ossessione sia lo studio di cose complesse, intrattabili ed incomprensibili.

A me stesso, per non aver abbandonato il progetto nelle settimane in cui il risultato principale sembrava impossibile da raggiungere e per aver finalmente capito che una scadenza è un teorema a cui si crede solo dopo averlo dimostrato.

Infine, a tutti quelli che ho dimenticato di nominare: l'omissione è un limite della mia memoria, ma non della mia gratitudine. Per assicurare la correttezza ho dovuto sacrificare la completezza, ma non posso neanche dimostrare la mia stessa correttezza. Dannato Gödel e i suoi due stupidi teoremi.

\noindent
\textit{Simone Bianco, Roma, Ottobre 2026}